% =============================================================================
% bab2-tinjauan-bisnis.tex — BAB II TINJAUAN PUSTAKA ASPEK BISNIS
% Proyek Bisnis / Studi Kasus — Lampiran II (S2)
% =============================================================================

\chapter{TINJAUAN PUSTAKA ASPEK BISNIS}
\label{bab:tinjauan-bisnis}

\section{Tinjauan Teori dan Konsep Bisnis}
\label{sec:tinjauan-teori-bisnis}

Tuliskan tinjauan teori dan konsep bisnis yang relevan dengan topik penelitian ini.
Uraikan kerangka teoretis yang digunakan sebagai landasan analisis.

\subsection{Teori/Konsep Pertama}
\label{subsec:konsep-pertama}

Uraikan teori atau konsep pertama yang relevan dengan permasalahan bisnis yang diteliti.
Berikan penjelasan mendalam mengenai teori ini beserta aplikasinya dalam konteks
bisnis \cite{contohSumber2025}.

\subsection{Teori/Konsep Kedua}
\label{subsec:konsep-kedua}

Uraikan teori atau konsep kedua yang relevan.

\subsection{Kerangka/Model Analisis Bisnis}
\label{subsec:model-analisis}

Jelaskan kerangka atau model analisis bisnis yang digunakan, misalnya:
\begin{itemize}
  \item Analisis SWOT (Strengths, Weaknesses, Opportunities, Threats)
  \item Business Model Canvas
  \item Porter's Five Forces
  \item Value Chain Analysis
  \item Dan lainnya sesuai topik penelitian
\end{itemize}

\section{Kajian Penelitian Terdahulu}
\label{sec:kajian-terdahulu-bisnis}

Berikut adalah ringkasan penelitian-penelitian terdahulu yang relevan dengan topik
penelitian ini:

\begin{table}[H]
\caption{Ringkasan Penelitian Terdahulu}
\label{tab:penelitian-terdahulu-pb}
\begin{isitable}
\begin{tabular}{p{0.5cm} p{3cm} p{4cm} p{3cm} p{3.5cm}}
\toprule
\textbf{No} & \textbf{Penulis (Tahun)} & \textbf{Judul/Topik} & \textbf{Metode/Alat Analisis} & \textbf{Temuan Utama} \\
\midrule
1 & Penulis (2023) & Topik penelitian & Metode analisis & Temuan utama \\
2 & Penulis (2022) & Topik penelitian & Metode analisis & Temuan utama \\
\bottomrule
\end{tabular}
\end{isitable}
\end{table}

\section{Kerangka Pemikiran}
\label{sec:kerangka-bisnis}

Berdasarkan kajian teori dan penelitian terdahulu yang telah diuraikan, kerangka
pemikiran dalam penelitian ini adalah sebagai berikut:

\begin{figure}[H]
\centering
% \includegraphics[width=0.85\textwidth]{figures/kerangka-bisnis}
\fbox{\parbox{12cm}{\centering\vspace{2cm}[Kerangka Pemikiran Proyek Bisnis]\vspace{2cm}}}
\caption{Kerangka Pemikiran Proyek Bisnis}
\label{fig:kerangka-bisnis}
\end{figure}
